\chapter{Theorieën van reactiesnelheden}\label{ch:b3:rate-theories}

In 2017 werd een geneesmiddel goedgekeurd dat alleen van een ouder middel
verschilt doordat de zes waterstofatomen van zijn twee methoxygroepen
deuteriumatomen zijn. De lever breekt het middel af door precies die
bindingen C--H te verbreken, en een binding C--D wordt meerdere keren
trager verbroken: het zware molecuul blijft langer in het bloed, en wordt
in kleinere, minder frequente doses ingenomen. Niets in de wet van
Arrhenius verklaart waarom een neutron in een kern een reactie zou
vertragen. Dit hoofdstuk leidt snelheidsconstanten af uit moleculaire
eigenschappen: eerst uit botsingen, daarna uit de vorm van het
\omterm{def:b3:computational-chemistry:pes}{potentiaalenergieoppervlak} en de toestandssommen van de
\omterm{def:b3:rate-theories:tst}{overgangstoestandstheorie}, die de grootte van zulke isotoopeffecten
voorspelt; het eindigt met reacties in oplossing, waar het oplosmiddel een
snelheidslimiet oplegt.

\begin{recall}
Het deel van jaar~1 definieerde de snelheidsconstante, de wet van
Arrhenius $k =
A\eu^{-E_a/RT}$ met haar activeringsenergie en pre-exponentiële factor,
elementaire stappen, de overgangstoestand, het energieprofiel langs de
reactiecoördinaat en het postulaat van Hammond. Het deel van jaar~2
definieerde ionsterkte en activiteitscoëfficiënten, met de limietwet van
Debye-Hückel $\log\gamma_i = -Az_i^2\sqrt I$ ($A \approx
0.51$ in water bij \qty{25}{\celsius}), en paste kleinstekwadratenrechten
aan.
\Cref{ch:b3:computational-chemistry} voerde potentiaalenergieoppervlakken
in;
\cref{ch:b3:partition-functions,ch:b3:statistical-thermo-applied}
berekenden evenwichtsconstanten uit toestandssommen en nulpuntsenergieën.
Uit de natuurkunde: de maxwell-boltzmannverdeling van de moleculaire
snelheden, en de wet van Fick voor diffusie.
\end{recall}

\begin{omfigure}
\includegraphics[width=0.6\linewidth]{images/book4/ai/b3-12-pharmacy.jpg}

{\small Geneesmiddelen op een apotheekschap. Waterstof vervangen door
deuterium op de bindingen die de lever het eerst aanvalt, kan de tijd
verlengen dat een middel in het lichaam blijft.}
\end{omfigure}

\section{Botsingstheorie}

Twee moleculen in een gas kunnen alleen reageren als ze elkaar ontmoeten.
Stel ze voor als harde bollen met diameters $d_{\mathrm A}$ en $d_{\mathrm B}$:
ze botsen wanneer hun middelpunten elkaar tot op $d =
\frac12(d_{\mathrm A} + d_{\mathrm B})$ naderen.

\begin{definition}[Botsingsdoorsnede, sterische factor]\label{def:b3:rate-theories:collision}
De \emph{botsingsdoorsnede}\index{botsingsdoorsnede} van twee moleculen
is de oppervlakte $\sigma = \pi d^2$ van de schijf, loodrecht op hun
relatieve snelheid, waarbinnen het middelpunt van het ene moet passeren om
het andere te raken. De
\emph{sterische factor}\index{sterische factor} $P$ is de verhouding van de
gemeten pre-exponentiële factor van een reactie tot de factor die de
botsingstheorie voorspelt.
\end{definition}

\begin{omfigure}
\begin{tikzpicture}[font=\small, >=Stealth]
  \shade[ball color=omDef!40] (0,0) circle (0.45);
  \node at (0,0) {A};
  \draw[black!50, dashed] (-0.0,0.85) -- (5.2,0.85) (-0.0,-0.85) -- (5.2,-0.85);
  \draw[black!60] (0,0) ellipse (0.25 and 0.85);
  \draw[<->] (-0.45,0) -- (-0.45,0.85) node[midway, left] {$d$};
  \shade[ball color=omThm!40] (3.3,0.55) circle (0.4);
  \node at (3.3,0.55) {B};
  \shade[ball color=omThm!20] (4.6,1.45) circle (0.4);
  \node at (4.6,1.45) {B};
  \draw[->, thick] (0.6,0) -- (2.2,0) node[midway, below] {$v_{\mathrm{rel}}\,\Delta t$};
  \node[anchor=west, align=left, font=\scriptsize] at (5.3,0) {cilinder met\\ doorsnede\\ $\sigma = \pi d^2$};
  \node[font=\scriptsize, anchor=south] at (3.3,1.0) {geraakt};
  \node[font=\scriptsize, anchor=south] at (4.6,1.9) {gemist};
\end{tikzpicture}

{\small De botsingscilinder. Molecuul A, dat beweegt met de relatieve
snelheid $v_{\mathrm{rel}}$, doorloopt
in een tijd $\Delta t$ een cilinder met doorsnede $\sigma = \pi d^2$, met
$d$ de som van de twee stralen; elk B waarvan het middelpunt erin ligt,
wordt geraakt.}
\end{omfigure}

\begin{theorem}[Snelheidsconstante volgens de botsingstheorie]\label{thm:b3:rate-theories:collision}
Als elke botsing waarvan de kinetische energie langs de verbindingslijn
van de middelpunten een drempel $E_0$ (per mol) overschrijdt, tot reactie
leidt, is de snelheidsconstante van $\mathrm{A + B} \to$
producten
\[
  k = \sigma\bar v_{\mathrm{rel}}N_A\,\eu^{-E_0/RT}, \qquad \bar v_{\mathrm{rel}} = \sqrt{\frac{8kT}{\pi\mu}},\
  \mu = \frac{m_{\mathrm A}m_{\mathrm B}}{m_{\mathrm A} + m_{\mathrm B}} .
\]
\end{theorem}

\begin{proof}[Gedeeltelijk bewijs]
In een tijd $\Delta t$ doorloopt een molecuul A een volume $\sigma v_{\mathrm{rel}}\Delta t$ en raakt het de
$n_{\mathrm B}\sigma v_{\mathrm{rel}}\Delta t$ moleculen B die erin zitten: de
botsingsfrequentie per volume-eenheid
is $\sigma\bar v_{\mathrm{rel}}n_{\mathrm A}n_{\mathrm B}$, met $\bar v_{\mathrm{rel}}$
de gemiddelde relatieve snelheid die de
maxwell-boltzmannverdeling geeft (uit de natuurkunde overgenomen, met de
\omterm{def:b3:quantum-model-systems:oscillator}{gereduceerde massa}). Elke botsing wegen met haar snelheid en alleen die
behouden waarvan de energie langs de verbindingslijn van de middelpunten
de drempel overschrijdt, laat de fractie
$\eu^{-E_0/RT}$ van de botsingsfrequentie over (de integratie over de
botsingsparameter wordt aangenomen). Het omzetten van aantaldichtheden in
molaire concentraties vermenigvuldigt met
$N_A$.
\end{proof}

Omdat $\bar v_{\mathrm{rel}} \propto \sqrt T$, is de activeringsenergie $E_a = RT^2\,\dd\ln k/\dd T$ gelijk aan $E_0 + \frac12RT$,
dicht bij de drempel. De pre-exponentiële factor is van de orde van
$10^{11}$~\unit{L.mol^{-1}.s^{-1}} voor kleine moleculen. Gemeten factoren
zijn vaak veel kleiner: $P \ll 1$ voor moleculen die elkaar in een
bepaalde oriëntatie moeten ontmoeten, en de doorsnede heeft geen
eenduidige waarde voor echte moleculen, die geen harde bollen zijn. De
botsingstheorie geeft de grootteorde en de temperatuurafhankelijkheid;
ze kan $E_0$ niet geven, want daarvoor is het \omterm{def:b3:computational-chemistry:pes}{potentiaalenergieoppervlak}
nodig.

\section{Potentiaalenergieoppervlakken}

De reactie $\mathrm{H_A + H_BH_C \to H_AH_B + H_C}$ is de eenvoudigste die
er bestaat. Bij een collineaire nadering hangt de energie af van twee
afstanden, $r_{\mathrm{AB}}$ en $r_{\mathrm{BC}}$, en kan ze als een kaart
getekend worden.

\begin{definition}[Zadelpunt, pad van minimale energie]\label{def:b3:rate-theories:saddle}
Een \emph{zadelpunt}\index{zadelpunt} van een \omterm{def:b3:computational-chemistry:pes}{potentiaalenergieoppervlak} is
een punt waar de gradiënt nul is, en de energie maximaal is in één
richting en minimaal in alle andere. Het
\emph{pad van minimale energie}\index{pad van minimale energie}
is het pad van steilste daling vanaf het zadelpunt naar het dal van de
beginstoffen en dat van de producten; de lengte gemeten langs dat pad is
de reactiecoördinaat, en het zadelpunt is de overgangstoestand.
\end{definition}

\begin{omfigure}
\begin{tikzpicture}
\begin{axis}[ompes, width=0.62\linewidth, xmin=0.4, xmax=3.0, ymin=0.4, ymax=3.0,
  xlabel={$r_{\mathrm{AB}}$ / \AA}, ylabel={$r_{\mathrm{BC}}$ / \AA},
  colorbar, colorbar style={ylabel={$E$ / eV}, ylabel style={font=\small},
  yticklabel style={font=\small}}, point meta min=-4.7, point meta max=-0.8]
\addplot3[contour prepared={labels=false}, thick] table {figdata/out/bachelor-3-pes-collinear-contours.dat};
\addplot3[ompath] table[x=x, y=y, z expr=0] {figdata/out/bachelor-3-pes-collinear-mep.dat};
\node[omsaddle] at (axis cs:0.918,0.918,0) {};
\node[font=\small, anchor=south west] at (axis cs:0.95,0.95,0) {zadelpunt};
\node[font=\small, anchor=west, fill=white, fill opacity=0.85, text opacity=1, inner sep=1.5pt] at (axis cs:2.05,1.0,0) {$\mathrm{H_A + H_BH_C}$};
\node[font=\small, anchor=west, fill=white, fill opacity=0.85, text opacity=1, inner sep=1.5pt] at (axis cs:0.95,2.55,0) {$\mathrm{H_AH_B + H_C}$};
\node[font=\small, anchor=north east] at (axis cs:2.95,2.95,0) {$\ce{H + H + H}$};
\end{axis}
\end{tikzpicture}

{\small Een modelpotentiaalenergieoppervlak voor collineair $\ce{H + H2}$
(LEPS-vorm opgebouwd uit de morsekromme van \ce{H2}, met een
Sato-parameter van 0.17, gekozen als modelconstante). Hoogtelijnen van
\qty{-4.65}{eV} tot \qty{-0.8}{eV}, met als nulpunt van de energie drie
gescheiden atomen; de dalen liggen bij \qty{-4.75}{eV}, de diepte van de
put van \ce{H2}. Gestreept: het \omterm{def:b3:rate-theories:saddle}{pad van minimale energie}, dat het
symmetrische \omterm{def:b3:rate-theories:saddle}{zadelpunt} kruist bij $r_{\mathrm{AB}} = r_{\mathrm{BC}} = \qty{0.92}{\angstrom}$,
in dit model \qty{0.27}{eV} boven de dalen.}
\end{omfigure}

Het pad klimt uit het dal van de beginstoffen, waar $r_{\mathrm{BC}}$ dicht
bij de bindingslengte van \ce{H2} blijft terwijl A nadert, kruist het
\omterm{def:b3:rate-theories:saddle}{zadelpunt}, en daalt af in het dal van de producten. Bij het \omterm{def:b3:rate-theories:saddle}{zadelpunt} zijn
beide bindingen uitgerekt tot \qty{0.92}{\angstrom}: de oude binding is
half verbroken en de nieuwe half gevormd, en de energiekost van dit
compromis is de barrière. Voor de symmetrische reactie ligt het \omterm{def:b3:rate-theories:saddle}{zadelpunt}
op de diagonaal; voor een exotherme reactie schuift het naar het dal van
de beginstoffen (een vroege barrière, een overgangstoestand die op de
beginstoffen lijkt, zoals het postulaat van Hammond zegt), voor een
endotherme naar het dal van de producten (een late barrière).
De positie bepaalt welke energie helpt: translatie-energie, gericht langs
het ingangsdal, brengt de moleculen over een vroege barrière; een late
barrière, om een hoek heen, wordt gemakkelijker genomen met vibratie-energie
in de binding die verbroken wordt. Dit zijn de regels van Polanyi,
bevestigd door experimenten met moleculaire bundels.

\section{Overgangstoestandstheorie}

\begin{definition}[Geactiveerd complex]\label{def:b3:rate-theories:activated-complex}
Het \emph{geactiveerd complex}\index{geactiveerd complex} van een
elementaire stap is de verzameling configuraties van het reagerende
systeem in een dunne laag rond het \omterm{def:b3:rate-theories:saddle}{zadelpunt}, loodrecht op het \omterm{def:b3:rate-theories:saddle}{pad van
minimale energie}; het wordt behandeld als een deeltje $\ce{AB^{\ddagger}}$,
met een eigen toestandssom.
\end{definition}

\begin{definition}[Overgangstoestandstheorie, transmissiecoëfficiënt]\label{def:b3:rate-theories:tst}
De \emph{overgangstoestandstheorie}\index{overgangstoestandstheorie} neemt
aan dat de \omterm{def:b3:rate-theories:activated-complex}{geactiveerde complexen} in evenwicht zijn met de beginstoffen,
dat elk complex dat het gebied van het \omterm{def:b3:rate-theories:saddle}{zadelpunt} in de richting van de
producten doorkruist, die ook vormt, en dat de beweging langs de
reactiecoördinaat een vrije translatie is. De
\emph{transmissiecoëfficiënt}\index{transmissiecoefficient@transmissiecoëfficiënt} $\kappa \le 1$ corrigeert voor
de complexen die terugkeren naar de beginstoffen.
\end{definition}

\begin{theorem}[Vergelijking van Eyring]\label{thm:b3:rate-theories:eyring}
Voor een elementaire stap $\mathrm{A + B} \to \mathrm{AB}^{\ddagger} \to$ producten is
\[
  k = \kappa\frac{k_BT}{h}\overline{K}{}^{\ddagger}, \qquad
  \overline{K}{}^{\ddagger} = \frac{\bar q^{\ddagger}/N_A}{(q_{\mathrm A}/N_A)(q_{\mathrm B}/N_A)}\,\eu^{-\Delta E_0^{\ddagger}/RT},
\]
waarin $\bar q^{\ddagger}$ de toestandssom is van het \omterm{def:b3:rate-theories:activated-complex}{geactiveerde complex}
zonder de beweging langs de reactiecoördinaat, $q$ de molaire
toestandssommen per volume-eenheid zijn, en $\Delta E_0^{\ddagger}$ de
hoogte is van het nulpuntsniveau van het complex boven dat van de
beginstoffen. In thermodynamische vorm, met $c^\circ = \qty{1}{mol/L}$ en
de moleculariteit $m$,
\[
  k = \kappa\frac{k_BT}{h}(c^\circ)^{1-m}\,\eu^{\Delta^{\ddagger}S^\circ/R}\,\eu^{-\Delta^{\ddagger}H^\circ/RT}.
\]
Deze uitspraak is de \emph{vergelijking van Eyring}\index{vergelijking van Eyring}; $k_B$ is de
boltzmannconstante, zo geschreven om verwarring met $k$ te vermijden.
\end{theorem}

\begin{proof}
Quasi-evenwicht: $[\ce{AB^{\ddagger}}] = K^{\ddagger}[\mathrm A][\mathrm B]$, met $K^{\ddagger}$ uit de toestandssommen
(\cref{thm:b3:statistical-thermo-applied:k-from-q}, geschreven met
concentraties). Van de vibraties van het complex is er één de beweging
langs de reactiecoördinaat: een zeer slappe vibratie met frequentie $\nu \ll k_BT/h$, waarvan de
toestandssom $k_BT/h\nu$ is (de hogetemperatuurlimiet van
\cref{prop:b3:partition-functions:q-vib}); dus $q^{\ddagger} = \bar q^{\ddagger}k_BT/h\nu$. De
complexen steken naar de producten over met de frequentie $\nu$ van deze
beweging, zodat de
snelheid $\nu[\ce{AB^{\ddagger}}] = \nu\frac{k_BT}{h\nu}\overline{K}{}^{\ddagger}[\mathrm A][\mathrm B]$ is: de onbekende $\nu$ valt weg. Schrijven van
$\overline{K}{}^{\ddagger} = (c^\circ)^{1-m}\eu^{-\Delta^{\ddagger}G^\circ/RT}$ en $\Delta^{\ddagger}G^\circ = \Delta^{\ddagger}H^\circ - T\Delta^{\ddagger}S^\circ$ geeft de
tweede vorm, met $\kappa$ ingevoegd voor het terugkeren.
\end{proof}

De factor $k_BT/h$ is \qty{6.2e12}{s^{-1}} bij \qty{298}{K}: de snelheid van
een monomoleculaire stap zonder barrière en zonder entropieverlies. Alles
wat specifiek is voor de reactie, zit in de activeringsgrootheden.

\begin{definition}[Gibbsenergie, enthalpie en entropie van activering]\label{def:b3:rate-theories:activation}
De \emph{gibbsenergie van activering}\index{gibbsenergie van activering}
$\Delta^{\ddagger}G^\circ$, \emph{activeringsenthalpie}\index{activeringsenthalpie} $\Delta^{\ddagger}H^\circ$ en
\emph{activeringsentropie}\index{activeringsentropie} $\Delta^{\ddagger}S^\circ$
van een elementaire stap zijn de standaardgibbsenergie, -enthalpie en
-entropie van de vorming van het \omterm{def:b3:rate-theories:activated-complex}{geactiveerde complex} (zonder zijn
beweging langs de reactiecoördinaat) uit de beginstoffen, gedefinieerd
door de \omterm{thm:b3:rate-theories:eyring}{vergelijking van Eyring}.
\end{definition}

\begin{proposition}[Activeringsenergie en activeringsenthalpie]\label{prop:b3:rate-theories:ea-dh}
Voor een reactie in oplossing is $E_a = \Delta^{\ddagger}H^\circ + RT$; voor
een bimoleculaire gasreactie
$E_a = \Delta^{\ddagger}H^\circ + 2RT$.
\end{proposition}

\begin{proof}
$E_a = RT^2\,\dd\ln k/\dd T$, en $\ln k = \ln T + \Delta^{\ddagger}S^\circ/R - \Delta^{\ddagger}H^\circ/RT + $ const: $E_a =
RT + \Delta^{\ddagger}H^\circ$ wanneer de activeringsparameters niet met $T$
variëren. In oplossing is dat het resultaat. In de gasfase bevat de
betrekking van Van 't Hoff voor $K^{\ddagger}$, geschreven in concentraties,
de inwendige energie, $\Delta^{\ddagger}U^\circ = \Delta^{\ddagger}H^\circ - \Delta n^{\ddagger}RT$
met $\Delta n^{\ddagger} = 1 - m$: voor $m = 2$ is $E_a = \Delta^{\ddagger}H^\circ + 2RT$.
\end{proof}

De \omterm{def:b3:rate-theories:activation}{activeringsentropie} leest de structuur van de overgangstoestand. Twee
moleculen die tot één complex samengaan, verliezen translatie- en
rotatievrijheid: $\Delta^{\ddagger}S^\circ$ is sterk negatief, $-50$ tot
$-150$ \unit{J.K^{-1}.mol^{-1}}, voor associatieve stappen en cyclische
overgangstoestanden. Een molecuul dat op weg naar dissociatie een binding
losser maakt, wint vrijheid: $\Delta^{\ddagger}S^\circ > 0$.

\begin{proposition}[Standaardfouten van een aangepaste rechte]\label{prop:b3:rate-theories:slope-uncertainty}
Voor een kleinstekwadratenrechte $y = a + bx$ door $n$ punten waarvan de
$y$ onafhankelijke fouten met dezelfde variantie dragen, geschat door $s^2 = \sum_i(y_i - a - bx_i)^2/(n - 2)$,
hebben helling en afsnijding de standaardfouten
\[
  s_b = \frac{s}{\sqrt{\sum_i(x_i - \bar x)^2}}, \qquad
  s_a = s\sqrt{\frac1n + \frac{\bar x^2}{\sum_i(x_i - \bar x)^2}} .
\]
\end{proposition}

\begin{proof}
$b = \sum_i(x_i - \bar x)y_i/S_{xx}$, met $S_{xx} = \sum_i(x_i - \bar x)^2$,
is een lineaire combinatie van onafhankelijke $y_i$ met variantie $\sigma^2$: haar variantie is $\sigma^2\sum_i(x_i - \bar x)^2/S_{xx}^2 = \sigma^2/S_{xx}$.
Evenzo is $a = \bar y - b\bar x$, en $\bar y$ en $b$ zijn ongecorreleerd
(hun covariantie is
$\sigma^2\sum_i(x_i - \bar x)/(nS_{xx}) = 0$), zodat $\operatorname{var}a = \sigma^2/n + \bar x^2\sigma^2/S_{xx}$. Het vervangen van $\sigma^2$
door zijn schatting $s^2$ (de deler $n - 2$ houdt rekening met de twee
aangepaste parameters, hier aangenomen) geeft het resultaat.
\end{proof}

\begin{method}[Een eyringgrafiek met onzekerheden]\label{met:b3:rate-theories:eyring-plot}
\begin{enumerate}
  \item Meet $k$ bij vijf of meer temperaturen, gespreid over 30 tot \qty{40}{K}.
  \item Zet $y = \ln(k/T)$ uit tegen $x = 1/T$ en pas de kleinstekwadratenrechte $y = a + bx$ aan.
  \item $\Delta^{\ddagger}H^\circ = -Rb$ en $\Delta^{\ddagger}S^\circ = R[a - \ln(k_B/h)]$ ($\ln(k_B/h) = 23.760$ met $k$ in
        \unit{s^{-1}}).
  \item Hun standaardfouten zijn $Rs_b$ en $Rs_a$; een
        betrouwbaarheidsinterval van 95\,\% vermenigvuldigt ze met de $t$
        van Student voor $n - 2$ vrijheidsgraden (3.18 voor vijf punten).
  \item De afsnijding ligt ver buiten het gemeten bereik, bij $1/T = 0$: $\Delta^{\ddagger}S^\circ$
        is altijd veel minder nauwkeurig dan $\Delta^{\ddagger}H^\circ$, en de
        twee fouten zijn gecorreleerd.
\end{enumerate}
\end{method}

\begin{omfigure}
\begin{tikzpicture}
\begin{axis}[width=0.74\linewidth, height=6.2cm, xmin=3.12, xmax=3.62, ymin=-12.8, ymax=-6.4,
  axis lines=left, xlabel={$1000\,\mathrm K/T$}, ylabel={$\ln\bigl(k/(\mathrm s^{-1}\,\mathrm K^{-1})\bigr)$},
  tick label style={font=\small}, label style={font=\small},
  legend style={font=\scriptsize, draw=none, fill=white, fill opacity=0.85, text opacity=1,
  at={(0.98,0.98)}, anchor=north east}, legend cell align=left]
\addplot[omDef, thick] table[x=invT, y=fit] {figdata/out/bachelor-3-eyring-band-H.dat};
\addplot[omThm, thick] table[x=invT, y=fit] {figdata/out/bachelor-3-eyring-band-D.dat};
\addplot[omDef, thin, dashed, forget plot] table[x=invT, y=lo] {figdata/out/bachelor-3-eyring-band-H.dat};
\addplot[omDef, thin, dashed, forget plot] table[x=invT, y=hi] {figdata/out/bachelor-3-eyring-band-H.dat};
\addplot[omThm, thin, dashed, forget plot] table[x=invT, y=lo] {figdata/out/bachelor-3-eyring-band-D.dat};
\addplot[omThm, thin, dashed, forget plot] table[x=invT, y=hi] {figdata/out/bachelor-3-eyring-band-D.dat};
\addplot[only marks, mark=*, mark size=1.8pt, omDef, forget plot] table[x=invT, y=lnkT_H] {figdata/out/bachelor-3-eyring-points.dat};
\addplot[only marks, mark=square*, mark size=1.7pt, omThm, forget plot] table[x=invT, y=lnkT_D] {figdata/out/bachelor-3-eyring-points.dat};
\legend{middel met \ce{OCH3}, middel met \ce{OCD3}}
\end{axis}
\end{tikzpicture}

{\small Eyringgrafieken van de snelheidsconstanten van het weekendprobleem
(gegevens van het probleem) voor het geneesmiddel en zijn gedeutereerde
analoog, met de kleinstekwadratenrechten en hun
betrouwbaarheidsbanden van 95\,\% (gestreept, nauwelijks breder dan de
rechten: de punten spreiden ongeveer 1\,\%). De rechten zijn evenwijdig
binnen hun onzekerheid: dezelfde \omterm{def:b3:rate-theories:activation}{activeringsentropie}, activeringsenthalpieën
die ongeveer \qty{5}{kJ/mol} verschillen.}
\end{omfigure}

\section{Kinetische isotoopeffecten}

\begin{definition}[Kinetische isotoopeffecten]\label{def:b3:rate-theories:kie}
Het \emph{kinetisch isotoopeffect}\index{kinetisch isotoopeffect} van een
stap is de verhouding $k_{\mathrm{licht}}/k_{\mathrm{zwaar}}$ van zijn
snelheidsconstanten voor twee \omterm{def:b3:rovibrational-spectroscopy:rotational-constant}{isotopologen}, meestal
$k_{\mathrm H}/k_{\mathrm D}$. Het is een \emph{primair kinetisch
isotoopeffect}\index{primair kinetisch isotoopeffect}
wanneer de binding met het gesubstitueerde atoom in de stap gevormd of
verbroken wordt, een \emph{secundair kinetisch
isotoopeffect}\index{secundair kinetisch isotoopeffect}
wanneer dat niet zo is.
\end{definition}

\begin{proposition}[Maximaal primair kinetisch isotoopeffect]\label{prop:b3:rate-theories:primary-kie}
Als de strekvibratie van een binding X--H de reactiecoördinaat wordt,
zodat haar \omterm{def:b3:quantum-model-systems:oscillator}{nulpuntsenergie} in de overgangstoestand verloren gaat terwijl
alle andere bijdragen voor H en D dezelfde zijn, is
\[
  \frac{k_{\mathrm H}}{k_{\mathrm D}} = \exp\Bigl[\frac{hc(\tilde\nu_{\mathrm{XH}} - \tilde\nu_{\mathrm{XD}})}{2k_BT}\Bigr],
  \qquad \frac{\tilde\nu_{\mathrm{XH}}}{\tilde\nu_{\mathrm{XD}}} = \sqrt{\frac{\mu_{\mathrm{XD}}}{\mu_{\mathrm{XH}}}} .
\]
\end{proposition}

\begin{proof}
Het \omterm{def:b3:computational-chemistry:pes}{potentiaalenergieoppervlak} is voor beide \omterm{def:b3:rovibrational-spectroscopy:rotational-constant}{isotopologen} hetzelfde (de
elektronen zien het neutron niet). In de beginstof bevat de strekvibratie
X--H
$\frac12hc\tilde\nu_{\mathrm{XH}}$ aan \omterm{def:b3:quantum-model-systems:oscillator}{nulpuntsenergie}; in de
overgangstoestand is deze trilling de reactiecoördinaat geworden en bevat
ze niets. De barrière, gemeten tussen nulpuntsniveaus, is daardoor voor
H lager met $\frac12hc\tilde\nu_{\mathrm{XH}}$ en voor D met $\frac12hc\tilde\nu_{\mathrm{XD}}$: $\Delta E_0^{\ddagger}(\mathrm D) -
\Delta E_0^{\ddagger}(\mathrm H) = \frac12hc(\tilde\nu_{\mathrm{XH}} - \tilde\nu_{\mathrm{XD}})$, en de \omterm{thm:b3:rate-theories:eyring}{vergelijking van Eyring} geeft de
verhouding wanneer de andere factoren wegvallen. Het golfgetal van een
strekvibratie is evenredig met
$1/\sqrt\mu$ (\cref{ch:b3:rovibrational-spectroscopy}).
\end{proof}

\begin{omfigure}
\begin{tikzpicture}[font=\small, >=Stealth, x=1cm, y=1cm]
  \draw[->] (-1.4,0) -- (-1.4,4.7) node[anchor=south] {$E$};
  \draw[->] (-1.4,0) -- (10.6,0) node[anchor=north east] {reactiecoördinaat};
  \draw[very thick, black!70] plot[smooth, tension=0.6] coordinates {(1.6,3.4) (2.1,1.6) (2.7,0.8)
    (3.3,0.6) (3.9,0.8) (4.7,2.0) (5.5,3.6) (6.2,4.0) (6.9,3.6) (7.7,2.2) (8.5,1.2) (9.2,1.0)
    (9.9,1.2) (10.4,1.7)};
  \draw[black!50, dashed] (2.8,0.6) -- (3.8,0.6);
  \draw[omDef, thick] (2.42,1.35) -- (4.2,1.35);
  \draw[omThm, thick] (2.5,1.05) -- (4.05,1.05);
  \node[omDef, font=\scriptsize, anchor=south] at (3.3,1.37) {C--H};
  \node[omThm, font=\scriptsize, anchor=north] at (3.3,1.03) {C--D};
  \draw[black, thick] (5.7,4.0) -- (6.7,4.0);
  \draw[black!40, densely dotted] (0.0,4.0) -- (5.7,4.0);
  \draw[black!40, densely dotted] (0.7,1.35) -- (2.42,1.35) (0.0,1.05) -- (2.5,1.05);
  \draw[<->, omThm] (0.0,1.05) -- (0.0,4.0) node[pos=0.5, left, font=\scriptsize] {$E_0(\mathrm D)$};
  \draw[<->, omDef] (0.7,1.35) -- (0.7,4.0) node[pos=0.5, right, font=\scriptsize] {$E_0(\mathrm H)$};
  \node[anchor=south, align=center, font=\scriptsize] at (6.2,4.05) {overgangstoestand: de strekvibratie is\\ de reactiecoördinaat geworden};
\end{tikzpicture}

{\small Oorsprong van het primaire kinetische isotoopeffect. In de
beginstof bevat de strekvibratie C--H meer \omterm{def:b3:quantum-model-systems:oscillator}{nulpuntsenergie} dan de
strekvibratie C--D (niveaus boven de bodem van de put, gestreept); in de
overgangstoestand is deze vibratie de reactiecoördinaat geworden, en
bereiken beide \omterm{def:b3:rovibrational-spectroscopy:rotational-constant}{isotopologen} dezelfde top. De barrière is voor D groter, met
het verschil van de nulpuntsenergieën.}
\end{omfigure}

\begin{omfigure}
\begin{tikzpicture}
\begin{axis}[width=0.7\linewidth, height=5.6cm, xmin=200, xmax=800, ymin=1, ymax=50,
  ymode=log, axis lines=left, xlabel={$T$ / K}, ylabel={$k_{\mathrm H}/k_{\mathrm D}$ (maximum)},
  tick label style={font=\small}, label style={font=\small}, log ticks with fixed point,
  ytick={1,2,5,10,20,50}, legend style={font=\scriptsize, draw=none, at={(0.98,0.98)},
  anchor=north east}, legend cell align=left]
\addplot[omThm, thick] table[x=T, y=OH] {figdata/out/bachelor-3-kie.dat};
\addplot[black!60, thick] table[x=T, y=NH] {figdata/out/bachelor-3-kie.dat};
\addplot[omDef, thick] table[x=T, y=CH] {figdata/out/bachelor-3-kie.dat};
\legend{O--H (\qty{3657}{cm^{-1}}), N--H (\qty{3337}{cm^{-1}}), C--H (\qty{2917}{cm^{-1}})}
\end{axis}
\end{tikzpicture}

{\small Maximaal \omterm{def:b3:rate-theories:kie}{primair kinetisch isotoopeffect} als functie van de
temperatuur voor het verlies van een strekvibratie C--H, N--H of O--H
(\omterm{def:b3:rovibrational-spectroscopy:anharmonic}{fundamentele overgangen} van \ce{CH4}, \ce{NH3} en \ce{H2O}; de
golfgetallen voor deuterium uit de \omterm{def:b3:quantum-model-systems:oscillator}{gereduceerde massa}'s). Bij
\qty{298}{K} is het 6.5 voor C--H; het daalt naar 1 naarmate de
temperatuur stijgt.}
\end{omfigure}

De maximale waarde voor C--H bij kamertemperatuur, ongeveer 7, is een
ijkpunt. Waargenomen effecten van 2 tot 7 wijzen erop dat de binding C--H
in de snelheidsbepalende stap verbroken wordt; kleinere waarden erop dat
de overgangstoestand een deel van de strekvibratie behoudt
(asymmetrische, vroege of late waterstofoverdrachten), of dat de stap niet
de traagste is. Waarden die bij kamertemperatuur veel groter zijn dan 7,
kunnen niet van nulpuntsenergieën komen: ze verraden \omterm{def:b3:quantum-model-systems:tunnelling}{tunneling}
(\cref{ch:b3:quantum-model-systems}), die een deeltje met de dubbele massa
veel minder vertoont. Secundaire effecten, afkomstig van veranderingen van
de buigvibraties van bindingen C--H naast het reactiecentrum, zijn klein,
meestal 0.8 tot 1.2 per deuteriumatoom; ze onderscheiden bijvoorbeeld een
koolstofatoom dat trigonaal wordt (sp$^3$ naar sp$^2$, $k_{\mathrm H}/k_{\mathrm D} > 1$) van een dat
tetraëdrisch wordt ($< 1$).

\begin{method}[Een kinetisch isotoopeffect meten door competitie]\label{met:b3:rate-theories:kie}
\begin{enumerate}
  \item Voer de reactie uit op een mengsel van de twee \omterm{def:b3:rovibrational-spectroscopy:rotational-constant}{isotopologen}, of op
        een molecuul met H op één plaats en D op een equivalente plaats.
  \item Stop bij een lage omzetting en meet de verhouding van de producten
        met H en met D door massaspectrometrie of NMR.
  \item De productverhouding, gecorrigeerd voor de beginverhouding, is
        $k_{\mathrm H}/k_{\mathrm D}$; beide \omterm{def:b3:rovibrational-spectroscopy:rotational-constant}{isotopologen} ondergaan precies
        dezelfde omstandigheden, wat de methode nauwkeuriger maakt dan twee
        afzonderlijke snelheidsmetingen.
\end{enumerate}
\end{method}

\section{Reacties in oplossing}

In een vloeistof vliegt een molecuul tussen twee botsingen niet vrij rond:
het rammelt in een kooi van buren, waartegen het vele keren botst voordat
het wegdiffundeert. Twee beginstoffen die elkaar ontmoeten, blijven
gedurende vele botsingen samen: een ontmoeting.

\begin{definition}[Diffusie- en activeringscontrole, kooieffect]\label{def:b3:rate-theories:diffusion}
Een reactie in oplossing is
\emph{diffusiegecontroleerd}\index{diffusiegecontroleerde reactie}
wanneer de reactie bij elke ontmoeting snel is, zodat haar snelheid de
snelheid is waarmee de beginstoffen naar elkaar toe diffunderen; ze is
\emph{activeringsgecontroleerd}\index{activeringsgecontroleerde reactie}
wanneer slechts een kleine fractie van de ontmoetingen tot reactie leidt.
Het \emph{kooieffect}\index{kooieffect}
is de opsluiting van een paar moleculen door het omringende oplosmiddel,
waardoor ze tijdens één ontmoeting herhaaldelijk botsen.
\end{definition}

\begin{proposition}[Limiet van Smoluchowski]\label{prop:b3:rate-theories:smoluchowski}
De snelheidsconstante van een \omterm{def:b3:rate-theories:diffusion}{diffusiegecontroleerde reactie} tussen bollen
die reageren op de afstand $R^*$ is $k_d = 4\pi N_A(D_{\mathrm A} + D_{\mathrm B})R^*$. Met de betrekking van
Stokes-Einstein
$D = k_BT/6\pi\eta a$ voor bollen met straal $a$ in een oplosmiddel met
viscositeit $\eta$, en $R^* = a_{\mathrm A} + a_{\mathrm B}$ met $a_{\mathrm A} = a_{\mathrm B}$, is
\[
  k_d = \frac{8RT}{3\eta} .
\]
\end{proposition}

\begin{proof}[Gedeeltelijk bewijs]
Leg A vast in de oorsprong; B diffundeert met de relatieve
diffusiecoëfficiënt $D =
D_{\mathrm A} + D_{\mathrm B}$ en wordt vernietigd bij $r = R^*$. In de
stationaire toestand is de radiale flux door elke bol dezelfde, $J = 4\pi r^2D\,\dd c/\dd r$ (wet van
Fick, uit de natuurkunde); integreren van $R^*$, waar $c = 0$, tot
oneindig, waar $c = c_{\mathrm B}$, geeft
$J = 4\pi DR^*c_{\mathrm B}$, de snelheid per molecuul A. Met
Stokes-Einstein (aangenomen) is
$(D_{\mathrm A} + D_{\mathrm B})R^* = \frac{k_BT}{6\pi\eta}(\frac{1}{a} + \frac1a)(2a) = \frac{2k_BT}{3\pi\eta}$, en $k_d = 4\pi N_A \times
\frac{2k_BT}{3\pi\eta} = 8RT/3\eta$.
\end{proof}

\begin{example}[Snelheidslimieten in water en in hexaan]\label{ex:b3:rate-theories:speed-limit}
% ledger: visc:H2O.25C, visc:hexane.25C
Bij \qty{298.15}{K} geeft water ($\eta = \qty{0.890}{mPa.s}$) $k_d = 8 \times 8.314 \times
298.15/(3 \times \num{0.890e-3}) = \qty{7.4e6}{m^3.mol^{-1}.s^{-1}} = \qty{7.4e9}{L.mol^{-1}.s^{-1}}$; hexaan
($\eta = \qty{0.298}{mPa.s}$) \qty{2.2e10}{L.mol^{-1}.s^{-1}}. Geen enkele
bimoleculaire reactie tussen neutrale moleculen van gewone grootte is
sneller; gemeten constanten in de buurt van deze waarden (het uitdoven
van aangeslagen toestanden, de recombinatie van radicalen) zijn
\omterm{def:b3:rate-theories:diffusion}{diffusiegecontroleerd}. Ionen met tegengestelde lading trekken elkaar aan
en kunnen sneller gaan.
\end{example}

Ionen in oplossing voegen een elektrostatische bijdrage toe aan de
activering: de overgangstoestand van twee ionen heeft lading $z_{\mathrm A} + z_{\mathrm B}$, en zijn
activiteitscoëfficiënt verschilt van het product van de hunne.

\begin{definition}[Kinetisch zouteffect]\label{def:b3:rate-theories:salt}
Het \emph{kinetisch zouteffect}\index{kinetisch zouteffect} is de
verandering van de snelheidsconstante van een reactie tussen ionen met de
ionsterkte van de oplossing.
\end{definition}

\begin{proposition}[Vergelijking van Brønsted-Bjerrum]\label{prop:b3:rate-theories:bronsted-bjerrum}
Voor een stap tussen ionen met ladingen $z_{\mathrm A}$ en $z_{\mathrm B}$ in
verdunde waterige oplossing is
\[
  \log k = \log k_0 + 2Az_{\mathrm A}z_{\mathrm B}\sqrt I,
\]
waarin $k_0$ de snelheidsconstante bij ionsterkte nul is.
\end{proposition}

\begin{proof}
In de theorie van de overgangstoestand wordt het evenwicht met het \omterm{def:b3:rate-theories:activated-complex}{geactiveerde
complex} geschreven met activiteiten: $[\ce{AB^{\ddagger}}] = K^{\ddagger}[\mathrm A][\mathrm B]\gamma_{\mathrm A}\gamma_{\mathrm B}/\gamma^{\ddagger}$, dus $k = k_0\gamma_{\mathrm A}\gamma_{\mathrm B}/\gamma^{\ddagger}$.
Met de limietwet: $\log(\gamma_{\mathrm A}\gamma_{\mathrm B}/\gamma^{\ddagger}) = -A\sqrt I[z_{\mathrm A}^2 + z_{\mathrm B}^2 - (z_{\mathrm A} + z_{\mathrm B})^2] =
2Az_{\mathrm A}z_{\mathrm B}\sqrt I$.
\end{proof}

Ionen met hetzelfde teken reageren sneller wanneer er zout toegevoegd
wordt (de ionenatmosfeer schermt hun afstoting af), ionen met
tegengestelde tekens trager, en een neutrale beginstof vertoont geen
primair zouteffect: een grafiek van $\log k$ tegen $\sqrt I$ bij lage
ionsterkte meet het product van de ladingen van de reagerende deeltjes.

\begin{method}[Wat activeringsparameters over een mechanisme zeggen]\label{met:b3:rate-theories:mechanism-evidence}
\begin{enumerate}
  \item $\Delta^{\ddagger}S^\circ$ sterk negatief: een associatieve stap,
        een geordende of cyclische overgangstoestand, of een die het
        oplosmiddel ordent (ontstaan van lading); positief: een
        dissociatieve stap.
  \item Een primaire $k_{\mathrm H}/k_{\mathrm D}$ van 2 tot 7: de binding
        met H wordt in de snelheidsbepalende stap verbroken; boven ongeveer
        10 bij kamertemperatuur: \omterm{def:b3:quantum-model-systems:tunnelling}{tunneling}.
  \item De helling van $\log k$ tegen $\sqrt I$: het product $z_{\mathrm A}z_{\mathrm B}$ van de
        ladingen die elkaar in de snelheidsbepalende stap ontmoeten.
  \item Een snelheidsconstante rond $10^{10}$~\unit{L.mol^{-1}.s^{-1}} die varieert als $T/\eta$:
        diffusiecontrole.
\end{enumerate}
\end{method}

\begin{inthelab}[Een dubbelwandige cel voor een eyringstudie]
De reactie wordt gevolgd via haar absorbantie in een cel van een
spectrofotometer waarvan de mantel gevoed wordt door een thermostaatbad;
een thermokoppel in een referentiecel controleert de temperatuur tot op
\qty{0.1}{K}. Vijf tot zeven temperaturen, elk in drievoud gemeten, geven
een eyringgrafiek; de oplossingen worden vóór het mengen in de houder op
temperatuur gebracht, want een reactie die bij de verkeerde temperatuur
start, vertekent de eerste punten.
\end{inthelab}

\begin{history}[Eyring, Evans en Polanyi, 1935]
In 1935 publiceerden Henry Eyring, in Princeton, en Meredith Gwynne Evans
en Michael Polanyi, in Manchester, onafhankelijk van elkaar de theorie van
het \omterm{def:b3:rate-theories:activated-complex}{geactiveerde complex}. Eyring had in 1931 met Polanyi in Berlijn het
eerste \omterm{def:b3:computational-chemistry:pes}{potentiaalenergieoppervlak} voor $\ce{H + H2}$ berekend, met een
semi-empirische methode van het soort dat voor de figuur van dit hoofdstuk
gebruikt werd. De theorie werd met argwaan ontvangen: een tijdschrift
wees het artikel van Eyring eerst af als ondeugdelijk, en het werd pas
gepubliceerd nadat andere natuurkundigen ervoor ingestaan hadden.
\end{history}

\section{Oefeningen}

\begin{exercise}[$\star$]\label{exo:b3:rate-theories:1}
% ledger: kin:NO+O3.A
Neem voor $\ce{NO + O3 -> NO2 + O2}$ $\sigma = \qty{0.43}{nm^2}$ (gegevens van de oefening). Bereken $\bar v_{\mathrm{rel}}$
bij \qty{298}{K} ($\mu = \qty{18.46}{u}$), de pre-exponentiële factor
volgens de botsingstheorie, en de \omterm{def:b3:rate-theories:collision}{sterische factor}, gegeven de gemeten $A = \qty{2.07e-12}{cm^3.s^{-1}}$ per molecuul.
\end{exercise}

\begin{exercise}[$\star$]\label{exo:b3:rate-theories:2}
Een reactie in oplossing heeft $E_a = \qty{75.0}{kJ/mol}$ rond \qty{298}{K}.
Hoe groot is $\Delta^{\ddagger}H^\circ$? En voor een bimoleculaire
gasreactie met dezelfde $E_a$?
\end{exercise}

\begin{exercise}[$\star$]\label{exo:b3:rate-theories:3}
Voorspel het teken van $\Delta^{\ddagger}S^\circ$ voor: een
diels-aldercycloadditie; het monomoleculaire verlies van \ce{N2} uit een
azoverbinding; een $\mathrm{S_N2}$-reactie tussen twee neutrale moleculen
die twee ionen vormt in een polair oplosmiddel.
\end{exercise}

\begin{exercise}[$\star$]\label{exo:b3:rate-theories:4}
% ledger: visc:H2O.25C, visc:hexane.25C
Bereken de diffusiegecontroleerde snelheidsconstante in water en in hexaan
bij \qty{298.15}{K}. Welk oplosmiddel laat snellere ontmoetingen toe, en
waarom?
\end{exercise}

\begin{exercise}[$\star\star$]\label{exo:b3:rate-theories:5}
Een eersteordereactie heeft $k = \qty{1.00e-3}{s^{-1}}$ bij \qty{298.15}{K} en \qty{1.00e-2}{s^{-1}} bij
\qty{318.15}{K}. Bereken $\Delta^{\ddagger}H^\circ$ en $\Delta^{\ddagger}S^\circ$, en ter vergelijking $E_a$.
\end{exercise}

\begin{exercise}[$\star\star$]\label{exo:b3:rate-theories:6}
% ledger: vib:H2O.sym
Bereken het maximale primaire isotoopeffect voor de overdracht van een
proton O--H bij
\qty{298}{K} ($\tilde\nu_{\mathrm{OH}} = \qty{3657}{cm^{-1}}$, \omterm{def:b3:quantum-model-systems:oscillator}{gereduceerde massa}'s met O = 16, H = 1, D = 2).
\end{exercise}

\begin{exercise}[$\star\star$]\label{exo:b3:rate-theories:7}
% ledger: vib:CH4.nu1, vib:CD4.nu1
Een enzymatische waterstofoverdracht vanaf koolstof vertoont $k_{\mathrm H}/k_{\mathrm D} = 40$ bij \qty{298}{K}. Vergelijk
met het maximum uit nulpuntsenergieën ($\tilde\nu = 2917$ en \qty{2109}{cm^{-1}}) en trek een besluit.
\end{exercise}

\begin{exercise}[$\star\star$]\label{exo:b3:rate-theories:8}
Voorspel het effect van een verhoging van de ionsterkte van 0 naar 0.010
op de snelheid van: de oxidatie van \ce{I-} door \ce{S2O8^{2-}}
(snelheidsbepalende stap tussen deze twee ionen); de reactie van
\ce{CH3COOC2H5} met \ce{OH-}; de reactie van $\ce{[Co(NH3)5Br]^{2+}}$ met \ce{OH-}.
Geef de factor waar er een is ($A = 0.51$).
\end{exercise}

\begin{exercise}[$\star\star$]\label{exo:b3:rate-theories:9}
De reactie $\ce{F + H2 -> HF + H}$ is sterk exotherm. Waar ligt haar
barrière op het \omterm{def:b3:computational-chemistry:pes}{potentiaalenergieoppervlak}? Welke vorm van energie,
toegevoerd aan de beginstoffen, bevordert de omgekeerde reactie het best?
\end{exercise}

\begin{exercise}[$\star\star\star$]\label{exo:b3:rate-theories:10}
Een kleinstekwadratenaanpassing van $\ln(k/T)$ tegen $1/T$ over vijf punten geeft $b = -7446$ K
met $s_b = \qty{31}{K}$ en $a = 16.528$ met $s_a = 0.103$. Geef $\Delta^{\ddagger}H^\circ$ en $\Delta^{\ddagger}S^\circ$ met
hun standaardfouten en betrouwbaarheidsintervallen van 95\,\% ($t = 3.18$).
\end{exercise}

\begin{exercise}[$\star\star\star$]\label{exo:b3:rate-theories:11}
Toon aan, vertrekkend van de \omterm{thm:b3:rate-theories:eyring}{vergelijking van Eyring}, dat $E_a = \Delta^{\ddagger}H^\circ + RT$ voor een
monomoleculaire stap en dat $A = \eu\,(k_BT/h)\eu^{\Delta^{\ddagger}S^\circ/R}$. Leid $A$ af voor $\Delta^{\ddagger}S^\circ = 0$ bij
\qty{298}{K}, en vergelijk met de $10^{13}$ tot $10^{14}$~\unit{s^{-1}} van
eenvoudige bindingsbreuken.
\end{exercise}

\begin{exercise}[$\star\star\star$]\label{exo:b3:rate-theories:12}
Pas de \omterm{def:b3:rate-theories:tst}{overgangstoestandstheorie} toe op twee structuurloze atomen A en B
waarvan het \omterm{def:b3:rate-theories:activated-complex}{geactiveerde complex} een diatomisch molecuul is met
bindingslengte $d$ (geen vibratie meer over na het wegnemen van de
reactiecoördinaat, rotatie met $I = \mu d^2$). Toon aan dat ze de
snelheidsconstante van de botsingstheorie geeft met $\sigma = \pi d^2$.
\end{exercise}

\section{Probleem: een zwaarder geneesmiddel}

\begin{problem}[{Weekendprobleem --- een zwaarder geneesmiddel:
nulpuntsenergieën van de strekvibraties C--H en C--D, het maximale
isotoopeffect bij lichaamstemperatuur, de halveringstijd van een
gedeutereerd geneesmiddel, en een eyringanalyse van beide verbindingen}]
\label{pb:b3:rate-theories:1}
% ledger: vib:CH4.nu1, vib:CD4.nu1, imass:H1, imass:H2, hist4:deut.2017
Een geneesmiddel wordt langs twee wegen uit het lichaam verwijderd: 75\,\%
van zijn klaring door de O-demethylering van zijn twee \ce{OCH3}-groepen,
waarin de verbreking van een binding C--H snelheidsbepalend is, en 25\,\%
langs wegen die deze bindingen niet raken; zijn halveringstijd is
\qty{5.0}{h} (gegevens van het probleem). Zijn analoog draagt twee
\ce{OCD3}-groepen. Modellen voor de strekvibraties C--H en C--D: de
symmetrische strekvibraties van \ce{CH4}
(\qty{2917}{cm^{-1}}) en \ce{CD4} (\qty{2109}{cm^{-1}}). Snelheidsconstanten
van de eerste orde van de demethyleringsstap, gemeten met een preparaat
van leverenzymen (gegevens van het probleem):
\begin{center}\small
\begin{tabular}{lccccc}
\hline
$T$ / K & 278.15 & 288.15 & 298.15 & 308.15 & 318.15 \\
$k_{\mathrm H}$ / \unit{s^{-1}} & 0.00989 & 0.0263 & 0.0632 & 0.150 & 0.328 \\
$k_{\mathrm D}$ / \unit{s^{-1}} & 0.00117 & 0.00333 & 0.00874 & 0.0219 & 0.0503 \\ \hline
\end{tabular}
\end{center}

\medskip\noindent\textbf{Deel I --- Nulpuntsenergieën.}
\begin{enumerate}
  \item Bereken de \omterm{def:b3:quantum-model-systems:oscillator}{gereduceerde massa}'s van C--H en C--D (C = 12, H = 1.00783,
        D = 2.01410).
  \item Voorspel het golfgetal van C--D uit dat van C--H en vergelijk met de
        gemeten \qty{2109}{cm^{-1}}.
  \item Bereken de nulpuntsenergieën van de twee strekvibraties en hun
        verschil, in \unit{cm^{-1}} en \unit{kJ/mol} (gemeten golfgetallen).
  \item Wat gebeurt er in de overgangstoestand met de strekvibratie van de
        binding die verbroken wordt?
  \item Leid het verschil tussen de barrières voor H en D af.
\end{enumerate}

\medskip\noindent\textbf{Deel II --- Het isotoopeffect.}
\begin{enumerate}[resume]
  \item Bereken de maximale $k_{\mathrm H}/k_{\mathrm D}$ bij \qty{37}{\celsius}.
  \item Bereken ze bij \qty{25}{\celsius}.
  \item Waarom heet dit een maximum?
  \item Dragen de twee andere deuteriumatomen van elke \ce{CD3}-groep veel bij?
  \item Wat zegt een gemeten waarde rond 7 bij \qty{25}{\celsius} over het
        mechanisme?
\end{enumerate}

\medskip\noindent\textbf{Deel III --- Halveringstijden.}
\begin{enumerate}[resume]
  \item Hoe hangt de halveringstijd af van de totale klaringsconstante?
  \item Schrijf de totale snelheidsconstante van het H-middel als de som
        van zijn twee wegen, in fracties.
  \item Met hoeveel wordt de demethyleringsweg vertraagd in het D-middel?
  \item Bereken $k_{\mathrm{totaal}}(\mathrm D)/k_{\mathrm{totaal}}(\mathrm H)$.
  \item Bereken de halveringstijd van het D-middel.
  \item Bereken de verhouding van de halveringstijden.
  \item Hoe groot zou ze zijn als de hele klaring via demethylering verliep?
  \item Wat verandert een langere halveringstijd voor de patiënt?
\end{enumerate}

\medskip\noindent\textbf{Deel IV --- Een eyringanalyse.}
\begin{enumerate}[resume]
  \item Welke grafiek van de gegevens is volgens de \omterm{thm:b3:rate-theories:eyring}{vergelijking van Eyring}
        lineair?
  \item Pas ze aan voor het H-middel: $\Delta^{\ddagger}H^\circ$ en $\Delta^{\ddagger}S^\circ$ met hun standaardfouten.
  \item Doe hetzelfde voor het D-middel.
  \item Bereken $\Delta\Delta^{\ddagger}H^\circ = \Delta^{\ddagger}H^\circ(\mathrm D) - \Delta^{\ddagger}H^\circ(\mathrm H)$ en de standaardfout ervan.
  \item Is het gelijk aan het verschil in \omterm{def:b3:quantum-model-systems:oscillator}{nulpuntsenergie} van vraag 3?
  \item Wat zeggen de twee activeringsentropieën?
  \item Formuleer het resultaat: de voorspelde verhouding van de
        halveringstijden van de twee geneesmiddelen bij \qty{37}{\celsius}.
\end{enumerate}
\end{problem}
